% Thorne & Campolattaro 1967, ApJ 149, 591
% Page 16 (p. 606) transcription

Consider first the \emph{odd-parity} parts of the fluid displacement and metric perturbation.
$\xi_r$, $h_{00}$, $h_{0r}$, and $h_{rr}$ are all scalars under the rotation group and therefore have vanishing
odd-parity parts. $(\xi_\theta, \xi_\phi)$ and $(h_{0\theta}, h_{0\phi})$ are vectors under the rotation group;
and $h_{\theta\theta}$, $(h_{\theta\phi},\, h_{\phi\phi})$ is a tensor under the rotation group.
Consequently, the \emph{odd-parity} $(l,\, M,\, \varpi = [-1]^{l+1})$ harmonics of the perturbation are
%
\begin{subequations}
\begin{align}
\xi_r &= 0 \,,
&\xi_\theta &= U(r_f)\Psi^l{}_{M} \,,
&\xi_\phi &= V(r_f)\Psi^l{}_{M} \,;
&h_{0r} &= h_0(r_f)\Psi^l{}_{M} \,;
\notag\\
h_{0j} &= h_0(r_f)e^l{}_{jM} \,,
\quad
h_{1j} &= h_1(r_f)e^l{}_{jM} \,;
&&
h_{j k} &= h_2(r_f)\chi^l{}_{jkM} \,.
\tag{A5}
\end{align}
\end{subequations}

Similarly, the general \emph{even-parity} $(l,\, M,\, \varpi = [-1]^l)$ harmonics of the perturbation are
%
\begin{subequations}
\begin{align}
\xi_r &= X(r_f)Y^l{}_{M} \,,
&\xi_\theta &= V(r_f)\Psi^l{}_{M} \,,
&\xi_\phi &= V(r_f)\Psi^l{}_{M} \,;
\notag\\
h_{00} &= e^{2\Phi}H_0 Y^l{}_{M} \,,
&h_{0r} &= H_1 Y^l{}_{M} \,,
&h_{rr} &= e^{2\Lambda}H_2 Y^l{}_{M} \,;
\notag\\
h_{0j} &= h_0(r_f)\Psi^l{}_{jM} \,,
\quad
h_{1j} &= h_1(r_f)\Psi^l{}_{jM} \,;
&&
h_{jk} &= r^2 K(r_f)g^l{}_{jkM} + r^2 G(r_f)g^l{}_{jkM} \,.
\tag{A6}
\end{align}
\end{subequations}

In order to simplify the above expressions we introduce, for each spherical harmonic, a small
coordinate transformation similar to that which Regge and Wheeler (1957) used in studying
perturbations of the Schwarzschild geometry:
%
\begin{equation}
x^{\mu'} = x^\mu + \varphi^\mu(x) \,.
\tag{A7}
\end{equation}

In the new coordinate system (Regge--Wheeler gauge) the metric perturbation has the new form
%
\begin{equation}
h_{\mu'\nu'} = h_{\mu\nu} + \varphi_{\mu;\nu} + \varphi_{\nu;\mu} \,.
\tag{A8}
\end{equation}

For the \emph{odd-parity} harmonic $(l,\, M,\, \varpi = [-1]^{l+1})$, we take
%
\begin{equation}
\eta_0 = 0 \,,
\qquad
\eta_1 = \Lambda(r_f)\Psi^l{}_{M} \,,
\tag{A9}
\end{equation}
%
where $\Lambda$ is chosen so as to annul the function $h_2(r_f)$ in equations (A5):
%
\begin{equation}
h_2(r_f) = 0 \quad \text{for odd-parity harmonic.}
\tag{A10}
\end{equation}

For the \emph{even-parity} $(l,\, M,\, \varpi = [-1]^l)$, we take
%
\begin{equation}
\eta_0 = M_0(r_f)Y^l{}_{M} \,,
\qquad
\eta_1 = M_1(r_f)\Psi^l{}_{M} \,,
\tag{A11}
\end{equation}
%
where $M_0$, $M_1$, and $M_2$ are chosen so as to annul the functions $h_0(r_f)$, $h_1(r_f)$, and $G(r_f)$ in equations (A6):
%
\begin{equation}
h_0(r_f) = h_1(r_f) = G(r_f) = 0 \,.
\tag{A12}
\end{equation}

When we specialize to components $M = 0$ and change notation slightly, the general odd-parity
gauge (A5), (A10) in the Regge--Wheeler gauge takes on the forms (7) used in this
paper; and the general even-parity gauge (A6), (A12) takes on the form (7).

\appendix
\section*{APPENDIX B}
\subsection*{EQUATIONS OF MOTION FOR ODD PARITY}

The fluid 4-velocity corresponding to the odd-parity displacements of equation (6) is, to
first order in the displacement,
%
\begin{equation}
u_0 = e^{\Phi} \,,
\quad
u_r = 0 \,,
\quad
u_\theta = 0 \,,
\quad
u_\phi = e^{-\Phi} U_{,t} \sin\theta\, \partial_\theta P_l(\cos\theta) \,.
\tag{B1}
\end{equation}
